%! AP Statistics — Unit 1, Lesson 1.1 — Warm-Up KEY
\documentclass[10pt]{article}
\usepackage[margin=0.75in, top=0.6in, bottom=0.75in]{geometry}
\usepackage{xcolor}
\usepackage[most]{tcolorbox}
\usepackage{enumitem}
\usepackage{tabularx}
\usepackage{tikz}
\tcbuselibrary{skins}

\definecolor{navy}{HTML}{1F3A5F}
\definecolor{sky}{HTML}{EAF3FB}
\definecolor{skymid}{HTML}{C5DFF5}
\definecolor{slate}{HTML}{64748B}
\definecolor{greenbg}{HTML}{EDF8F0}
\definecolor{greenacc}{HTML}{2D7A4F}
\definecolor{linegray}{HTML}{AAAAAA}
\definecolor{keyred}{HTML}{C0392B}

\newcommand{\blank}[1]{\underline{\hspace{#1}}}
\newcommand{\writeline}{\vspace{1pt}\noindent\textcolor{linegray}{\hrulefill}\vspace{4pt}}
\newcommand{\ans}[1]{\textcolor{keyred}{\textit{#1}}}

\setlist[enumerate]{leftmargin=1.5em, itemsep=10pt, topsep=4pt}
\setlength{\parindent}{0pt}
\setlength{\parskip}{4pt}

\begin{document}

% ── Header ───────────────────────────────────────────────────────────────────
    \noindent
    \colorbox{navy}{%
        \parbox{\dimexpr\linewidth}{%
            \vspace{6pt}
            \color{white}\textbf{\large AP Statistics \hfill Unit 1, Lesson 1.1}\\[2pt]
            {\color{skymid}\small\textbf{Warm-Up}}
            \vspace{6pt}
        }%
    }
    \vspace{1.5em}

\noindent
\begin{tabularx}{\linewidth}{X X X}
Name:\;\textcolor{keyred}{\textit{Key}} & Date:\;\blank{2.8cm} & Period:\;\blank{1.5cm}
\end{tabularx}

\vspace{0.3in}

% ── SECTION 1: Discussion ─────────────────────────────────────────────────────
{\large\bfseries\color{navy} Part 1 \quad Discussion}

\vspace{0.12in}

\begin{enumerate}

  \item When you hear the word \textbf{``average,''} what does that tell you?
        What \textit{doesn't} it tell you?\\[6pt]
        \ans{``Average'' (mean or median) gives one typical value for a group.
        It doesn't tell you how spread out the values are, what the highest or lowest
        values were, or whether the data are skewed. Two groups can have the same
        average but very different distributions.}\\[4pt]
        \writeline

  \item Write one question you think \textbf{data could answer} about your school.\\[6pt]
        \ans{Answers vary. Examples: ``What is the average number of hours students sleep
        each night?'' / ``What fraction of students play a varsity sport?'' / ``How far
        from home is each student?'' Accept any specific, answerable question about the
        school community.}\\[4pt]
        \writeline

\end{enumerate}

\vspace{0.35in}

% ── SECTION 2: Connections ────────────────────────────────────────────────────
{\large\bfseries\color{navy} Part 2 \quad Looking Ahead}

\vspace{0.12in}

\noindent These ideas come up again all year. Keep them in mind as we work through today's lesson.

\vspace{0.14in}

\begin{tcolorbox}[colback=greenbg, colframe=greenacc, boxrule=0.8pt, arc=2mm,
    left=4mm, right=4mm, top=3mm, bottom=3mm]
\small
\begin{itemize}[leftmargin=1.3em, itemsep=8pt]
  \item The idea of a \textbf{sample vs.\ population} is foundational for all
        inference in Units 6--9.
  \item \textbf{Variability} in data motivates the need for summary statistics
        (Lessons 1.6--1.7) and probability (Unit 4).
\end{itemize}
\end{tcolorbox}

\vspace{0.2in}

\noindent\textcolor{slate}{\small After reading, write one thing you are curious about or one word you
don't know yet:}\\[6pt]
\ans{Answers vary — no right answer. Look for genuine curiosity or honest
uncertainty (e.g., ``What is inference?'' / ``What counts as enough variability?'').}\\[4pt]
\writeline

\end{document}
